\originalchapter{静电势、唯一性与电容}
\originalsection{泊松方程与边界条件}
静电场满足 $\nabla\times\ee=0$ 及 $\nabla\cdot\ee=\rho/\eps$. 在单连通区域可写 $\ee=-\nabla V$, 因而 $\nabla^2V=-\rho/\eps$. 对局域电荷、以无穷远为零势,
\[
V(\rr)=\frac1{4\pi\eps}\int\frac{\rho(\rr')}{|\rr-\rr'|}\dd^3x'.
\]
该式由命题\ref{prop:point}作用于积分核得到. 它适于已知电荷的开放空间问题; 对导体边界, 表面电荷通常未知, 必须求满足边界条件的势.

导体静电平衡时内部场为零, 故每个连通导体等势. 若法线由区域1指向区域2, 小药盒和小矩形环路给出 $\eps(E_{2n}-E_{1n})=\sigma$ 与 $\ee_{2t}=\ee_{1t}$. 导体外侧因而有 $E_n=\sigma/\eps$. 表面不能用自身全部外侧场计算受力: 薄层所受有效场是两侧平均场, 于是真空导体表面压力 $p=\sigma^2/(2\eps)$.

\begin{theorem}[狄利克雷唯一性]
有界连通区域内给定电荷密度和边界电势, 充分光滑的泊松方程解至多一个. 对外部区域还需指定无穷远行为, 使下述边界项消失.
\end{theorem}
\begin{proof}
两解之差 $u$ 满足 $\nabla^2u=0$ 且边界上 $u=0$. 用 $\nabla\cdot(u\nabla u)=|\nabla u|^2+u\nabla^2u$ 积分, 得 $\int_\Omega|\nabla u|^2\dd^3x=0$. 因而 $\nabla u=0$, 连通性使 $u$ 为常数, 零边界使常数为零.
\end{proof}
唯一性是镜像法的依据. 找到一个满足源、边界及无穷远条件的解, 就确定了物理解; 无须先知道导体的实际表面电荷. 若指定的是法向导数, 还需通量相容条件, 且电势只确定到常数.

\originalsection{场能与电容矩阵}
逐次组装静态电荷时, 每对电荷相互作用只计一次, 得 $U=\frac12\int\rho V\dd^3x$. 将泊松方程代入并分部积分, 对边界项消失的全空间问题,
\[
U=\frac\eps2\int|\nabla V|^2\dd^3x=\frac\eps2\int E^2\dd^3x.
\]
点电荷自能在此模型中发散, 不能把无穷自能直接解释成有限相互作用能. 假定导体外没有给定自由体电荷或额外固定源, 以无穷远零势为参考. 对多个导体, 线性泊松问题给出 $Q_i=C_{ij}V_j$, 且 $U=\frac12 C_{ij}V_iV_j$.

设 $u_i$ 为导体 $i$ 上取1、其他导体及无穷远取0的调和函数. 对任意电势, $V=\sum_iV_iu_i$. 分部积分给出 $C_{ij}=\eps\int\nabla u_i\cdot\nabla u_j\dd^3x$, 因而 $C_{ij}=C_{ji}$. 这也解释电容矩阵的互易性不是对几何对称性的猜测.

\begin{example}
同轴圆柱电容器内外半径为 $a,b$, 长度 $\ell\gg b$, 内导体电荷为 $q$, 忽略端部效应. 求电容.
\end{example}
\begin{solution}
高斯柱面给 $E_s=q/(2\pi\eps\ell s)$, 方向由内向外. 内导体电势高于外导体, 差值为 $q\ln(b/a)/(2\pi\eps\ell)$, 故 $C=2\pi\eps\ell/\ln(b/a)$. 当间隙 $d=b-a\ll a$ 时 $\ln(b/a)\simeq d/a$, 得 $C\simeq\eps(2\pi a\ell)/d$, 与展开后的平行板一致.
\end{solution}
固定电荷时机械力 $F_x=-\partial_x[Q^2/(2C)]$. 固定电压时电源也做功: $\dd U_\mathrm{source}=-V\dd Q$. 电场与电源的合势为 $-CV^2/2$, 所以 $F_x=(V^2/2)\partial_xC$. 忽略电源会把介质吸入力或极板吸引力的符号算反.

\originalsection{边界受力与能量的计算步骤}
在表面场跃变处, “把电荷乘上电场”先要回答电场是否包含这片电荷自身的场. 取局部近似平面的薄带电层, 两侧的自身场为 $\pm\sigma\nn/(2\eps)$; 来自其余电荷的场在层的两侧相同. 将两侧总场平均, 自身场抵消, 得 $\ee_\mathrm{other}=(\ee_\mathrm{out}+\ee_\mathrm{in})/2$. 因而面力密度为 $\sigma\ee_\mathrm{other}$. 对导体 $\ee_\mathrm{in}=0$, 真空外场 $\ee_\mathrm{out}=\sigma\nn/\eps$, 压力才是 $\sigma^2/(2\eps)$. 这个论证使用薄层极限和局部平面近似; 尖角处应先求可积的总力, 不强行给有限点值压力.

场能中的 $1/2$ 也应从建立过程判断. 把最终密度 $\rho$ 缓慢从零增至 $\lambda\rho$, $0\le\lambda\le1$. 线性真空理论使已有电势为 $\lambda V$, 新添电荷为 $\rho\dd\lambda$. 因而
\[
 U=\int_0^1\dd\lambda\int\lambda V\rho\dd^3x
 =\frac12\int\rho V\dd^3x.
\]
用 $\rho=-\eps\nabla^2V$ 后, 格林第一恒等式给
$\int V\nabla^2V=\oint V\partial_nV-\int|\nabla V|^2$.
对局域有限能量分布, 无穷远项消失. 若求的是一个有限区域而边界有非零势, 该边界项一般不能删. 例如电容器场能可以只在有场区域积分, 但 $\frac12\int\rho V$ 的 $\rho$ 要包含导体面的电荷.

固定电压的虚功可逐步算出. $Q=CV$ 给 $\dd Q=V\dd C$, 场能变化为 $\dd U_\mathrm{field}=V^2\dd C/2$, 电源损失为 $\dd U_\mathrm{source}=-V\dd Q=-V^2\dd C$. 总势能下降 $-V^2\dd C/2$, 故沿使 $C$ 增加的方向有正机械力. 固定电荷时没有电源交换, $\dd U=-Q^2\dd C/(2C^2)$, 所得方向相同. 本科作业2题4和作业6题6分别用这两个结论; 读者应先判定保持的是 $Q$ 还是 $V$, 再微分.

\originalsection{法向边界的相容条件}
若给定 $\partial_nV=h$ 而非边界势, 对泊松方程积分得到必要条件
$\oint_{\partial\Omega}h\dd a=-\eps^{-1}\int_\Omega\rho\dd^3x$. 它说指定向外的通量不能与总电荷矛盾. 两个同源、同法向数据解之差 $u$ 有 $\partial_nu=0$, 同一个能量积分仍给 $\nabla u=0$, 但无法再把常数定为零. 指定一个参考点电势或平均值即可除去该自由度. 这解释了为何“只求电场”与“求具有指定零点的势”在边界信息上不同.


本章依据 8.07 第3、4讲、作业2、3; 场能与边界适用条件在中文重编中一并说明.
